% ============================================================
% Section 6: Conclusion
% ============================================================

\section{Conclusion}
\label{sec:conclusion}

We present \textbf{\ours{}}, the first topology-preserving synthetic P\&ID
dataset, and demonstrate that training a patch-based Relationformer
pipeline exclusively on \ours{} achieves 63.8\% edge mAP on the
real-world OPEN100 benchmark with zero annotated training P\&IDs.
This is 30~percentage points above template-based synthetic training and
18~points above the prior modular baseline that requires real data.
A controlled comparison confirms that this gain is attributable entirely
to generation quality: the same model trained on random-placement
synthetic data achieves only ${\sim}33\%$, consistent with independent
findings in~\cite{sturmer2025}.
The scaling analysis identifies seed diversity as the primary constraint,
with clear diminishing returns beyond 400 synthetic images from 12 seeds.

The core practical message is simple: a company that can provide 12
real P\&IDs as seeds --- without sharing them externally --- can now
bootstrap a competitive P\&ID digitizer with no additional annotation.
Adding even a small number of real training images (E3~Mixed) closes
the remaining gap to within 2.8~pp of published SOTA.
\ours{}, the generation pipeline, and trained models will be publicly
released upon acceptance to enable reproducible research in the
data-scarce P\&ID digitization regime.
